\section{Networks of Authority, Labor, and Disagreement}
\label{sec:origen-relationships}

Origen never occupied a single institution that could authorize, finance, and
preserve all his work. Demetrius could entrust Alexandrian teaching and later
repudiate the Palestinian ordination; Alexander and Theoctistus could receive
the same man as preacher and presbyter. Ambrose could finance dictation but not
settle episcopal standing. Caesarean custody could preserve letters that the
Alexandrian proceedings did not. The life moved through overlapping powers
rather than one continuous office.

That structure explains why relationship terms must remain exact. Leonides was
father and teacher; Heraclas pupil, colleague, and successor; Gregory a
panegyrist-pupil; Ambrose a patron, prompter, and fellow confessor. Julius
Africanus was an erudite correspondent whose surviving challenge elicited a
substantial reply; the exchange does not prove personal friendship. Heracleon
and Celsus survive largely inside arguments written to defeat them. Each
relation opened a different kind of evidence, and none gives transparent access
to private motive.

The rupture with Alexandria displays the limits most sharply. Eusebius assigns
Demetrius envy; later Photius, \emph{Bibliotheca} codex 118, reports two
assemblies through Pamphilus's lost apology. No minutes, complete charges, or
votes survive. The evidence establishes divided authorization and lasting
separation, not a secure single motive or a modern trial transcript.

The manuscript workshop likewise distributes agency without dissolving
authorship. Eusebius 6.23.1--2 names more than seven alternating shorthand
writers, ordinary copyists, and women trained in fine writing. Ambrose provided
resources and questions; Origen dictated, argued, corrected, and attached his
name; workers made the output reproducible. The sources do not disclose wages,
individual assignments, or how far patronal pressure shaped particular
conclusions. They do establish that the celebrated corpus was material and
collaborative.

Access was unequal. Gregory's years of advanced study presupposed literacy,
travel, and leisure; congregational preaching reached a different public.
Named bishops and elite pupils dominate the archive, while a mother, a woman
host, Juliana, Herais, and women copyists enter only in bounded episodes. Their
presence corrects the solitary-genius picture, but the missing names do not
license a reconstructed common network.

Disagreement also produced preservation. Africanus's objection keeps both
sides of a textual argument in view. Origen's refutations conserve pieces of
Heracleon and Celsus while controlling their order and frame. Gregory's praise,
Porphyry's attack, Eusebius's defense, and the Alexandrian repudiation therefore
cannot be averaged into a neutral portrait; they must be read as acts by
participants with different objects and stakes.

The costs are correspondingly differentiated. Imperial authorities executed
Leonides and pupils and tortured Origen; Origen's own austerity endangered his
health; Alexandrian authorities ended his ministry there; later translators
and excerptors determined much of what readers could hear. These are not one
institutional mechanism and Origen did not cause every injury around him. They
show why his intellectual history is inseparable from bodies, money,
authorization, and the uneven custody of words.
